% The reference debrief: what /exercise writes at the end of a session.
% Specimen: no real session. Copy the structure, never the content.
\documentclass{kaelearn}
\subject{Analysis I}  \folio{Debrief}  \title{Heads and Tails}
\shorttitle{Specimen · 2026-10-16}
\subtitle{Two problems, one blind spot: who chooses $\varepsilon$.}
\problemparts
\begin{document}
\makeheader
\marginkael{proud}

\begin{panel}{The session}{read this first}{}
\begin{itemize}
  \item \textbf{Specimen 1}, a convergent sequence is bounded: solved yourself after rung 1.
  \item \textbf{Specimen 2}, $(-1)^n$ diverges: rung 3, then you finished it.
  \item \textbf{The pattern:} both times, $\varepsilon$ was yours to choose and you waited for the
        problem to hand you one. ``For every $\varepsilon$'' in a hypothesis is a promise you may
        cash in at any value; ``there is an $\varepsilon$'' in a negation is a choice you make.
        Choose the one that makes the rest easy.
  \item \textbf{Again on 23 October:} one variant per problem on the last page. They come back in
        your weekly review.
\end{itemize}
\end{panel}

\section{Convergent Means Bounded}

\begin{panel}[unbreakable]{The problem}{Specimen 1}{}
Let $(a_n)$ be a real sequence with $a_n\to a$. Show that $(a_n)$ is bounded: there is an $M$ with
$|a_n|\le M$ for all $n$.
\end{panel}

\subsection{Your path}
You started by trying to bound $|a_n|$ directly and stalled: there are infinitely many terms, and
nothing in the problem says how large any of them is. You wrote the definition of convergence down
correctly, then treated its $\varepsilon$ as something the problem had to supply. The question that
unlocked it (rung 1) was \emph{which $\varepsilon$ would you like?} You picked $\varepsilon=1$, saw
that every term past $N$ sits within $1$ of $a$, and finished the proof yourself.
\aside{The analysts of the 1800s needed decades to make ``eventually'' precise. You needed one
question.}

\begin{panel}[unbreakable]{Weak points}{named exactly}{}
\begin{enumerate}
  \item \textbf{You waited to be given $\varepsilon$.} In a hypothesis, ``for every $\varepsilon>0$
        there is an $N$'' lets you choose the $\varepsilon$. \emph{Fix:} the moment a convergence
        hypothesis appears, ask which $\varepsilon$ you would like.
  \item \textbf{You tried to handle every term in one move.} Infinitely many terms looked like
        infinitely many problems. \emph{Fix:} split the sequence into a finite head and a tail,
        and treat them separately.
\end{enumerate}
\end{panel}

\begin{panel}[unbreakable]{The deep idea}{head and tail}{}
Convergence says nothing about the first terms and everything about the tail. Finitely many numbers
are always harmless, because they have a maximum. So: tame the tail with one fixed $\varepsilon$,
then take the maximum over the head.
\end{panel}

\begin{center}
\begin{tikzpicture}
\begin{axis}[klplot,width=.88\linewidth,height=46mm,xmin=0.5,xmax=20.5,ymin=-1.4,ymax=6.4,
  xtick={1,5,10,15,20},ytick={0},xlabel={$n$},ylabel={$a_n$}]
  \fill[pattern={dither25},pattern color=kllow] (axis cs:6.5,1) rectangle (axis cs:20.5,3);
  \draw[kldim,dashed] (axis cs:0.5,2)--(axis cs:20.5,2) node[right,text=kldim]{$a$};
  \draw[kldim,dashed] (axis cs:0.5,5.4)--(axis cs:20.5,5.4) node[right,text=kldim]{$M$};
  \draw[kldim,dashed] (axis cs:6.5,-1.4)--(axis cs:6.5,6.4);
  \node[text=kldim,anchor=south west] at (axis cs:6.6,-1.4) {$N$};
  \node[text=kldim,font=\fell\itshape\small] at (axis cs:3.5,6.0) {head: take the max};
  \node[text=kldim,font=\fell\itshape\small] at (axis cs:13.5,3.5) {tail: within $1$ of $a$};
  \addplot[only marks,mark=*,mark size=1.4pt,klhi] coordinates {
    (1,5.4)(2,-0.6)(3,4.3)(4,0.4)(5,3.6)(6,1.1)(7,1.05)(8,2.83)(9,1.26)(10,2.67)(11,1.4)(12,2.55)
    (13,1.49)(14,2.48)(15,1.56)(16,2.42)(17,1.61)(18,2.37)(19,1.65)(20,2.33)};
\end{axis}
\end{tikzpicture}
\end{center}
\plotcaption{The head is six points: their maximum exists. The tail lives in the band, whose height
$\varepsilon=1$ fixed. $M$ covers both.}

\anim{heads-and-tails}{Watch $N$ while $\varepsilon$ shrinks: it runs right, the head grows, and
it never becomes infinite. 30 seconds.}

\subsection{Generalize it}
\begin{itemize}
  \item \textbf{Weaken convergence to Cauchy:} still bounded, by the same proof. Pick
        $\varepsilon=1$ and compare every tail term with $a_N$ instead of $a$.
  \item \textbf{Ask for the converse:} false. $(-1)^n$ is bounded and does not converge (Problem 02).
        Bounded \emph{and monotone} does converge, but that is the monotone convergence theorem,
        and it needs the completeness of $\mathbb{R}$.
  \item \textbf{Change the setting:} the split works for functions too. If $f$ is continuous on
        $[0,\infty)$ and $f(x)\to L$ as $x\to\infty$, then $f$ is bounded: $\varepsilon=1$ tames
        the tail $[R,\infty)$, and the head $[0,R]$ is compact, so $f$ has a maximum there.
\end{itemize}

\subsection{Same structure, elsewhere}
Every ``eventually'' statement in analysis has this shape: a property that holds from some $N$ on,
plus finitely many exceptions that a maximum or a minimum absorbs. You will meet it again when you
prove that Cauchy sequences are bounded, that a sequence with a nonzero limit stays away from $0$
(the variant at the end), and the product rule for limits, whose proof needs exactly this lemma.

\why{Why are you allowed to pick $\varepsilon=1$?}{Because convergence is your hypothesis here.
``For every $\varepsilon>0$ there is an $N$'' holds for every $\varepsilon$, so you may cash it in at
whatever value you like, and $1$ is the laziest choice. If you were \emph{proving} convergence, the
roles would flip: someone else picks $\varepsilon$, and you answer with an $N$.}

\why{Why is the head never the problem?}{A finite set of real numbers always has a largest element,
so $\max\{|a_1|,\dots,|a_{N-1}|\}$ exists. The argument needs the head to be finite. Without
convergence there may be no $N$ past which the terms stay in a band, and then the head is the whole
sequence: $a_n=n$.}

\subsection{The solution, idea first}
\emph{Idea:} tame the tail with $\varepsilon=1$, take the maximum over the head.
\begin{proof}
Apply the definition of $a_n\to a$ with $\varepsilon=1$: there is an $N$ with $|a_n-a|<1$ for all
$n\ge N$. For those $n$, the triangle inequality gives $|a_n|\le|a_n-a|+|a|<1+|a|$. The remaining
terms $a_1,\dots,a_{N-1}$ are finitely many, so
\[M=\max\{|a_1|,\dots,|a_{N-1}|,\,1+|a|\}\]
exists, and $|a_n|\le M$ for every $n$.
\end{proof}

\section{Divergence by Negation}

\begin{panel}[unbreakable]{The problem}{Specimen 2}{}
Show that the sequence $a_n=(-1)^n$ does not converge.
\end{panel}

\subsection{Your path}
You knew it diverges and first argued that it ``keeps jumping'', which is a picture, not a proof.
Your first negation of convergence read ``for every $N$ there is an $\varepsilon$'': the quantifiers
were swapped. Once they were in order, rung 3 gave the sub-goal: \emph{for any candidate limit $a$,
find one $\varepsilon$ that the sequence escapes infinitely often.} You chose $\varepsilon=1$ and
finished it.

\begin{panel}[unbreakable]{Weak points}{named exactly}{}
\begin{enumerate}
  \item \textbf{Negating nested quantifiers.} You swapped $\forall N$ and $\exists\varepsilon$.
        \emph{Fix:} negate mechanically from the left: every $\forall$ becomes $\exists$, every
        $\exists$ becomes $\forall$, the inequality flips, and the order stays.
  \item \textbf{Waiting for $\varepsilon$, again.} In the negation, ``there is an $\varepsilon$'' means
        you choose it. Same blind spot as Problem 01.
  \item \textbf{``Obviously'' as an argument.} \emph{Fix:} when something is obvious, write the
        definition it is obvious about, then say it in that language.
\end{enumerate}
\end{panel}

\begin{panel}[unbreakable]{The deep idea}{half the gap}{}
To show that a sequence diverges, find two values it visits infinitely often. Any candidate limit is
at least half their distance away from one of them, so $\varepsilon$ = half the gap defeats every
candidate at once.
\end{panel}

\begin{center}
\begin{tikzpicture}
\begin{axis}[klplot,width=.88\linewidth,height=40mm,xmin=0.5,xmax=14.5,ymin=-1.6,ymax=2,
  xtick={1,5,10},ytick={-1,0,1},xlabel={$n$},ylabel={$a_n$}]
  \fill[pattern={dither25},pattern color=kllow] (axis cs:0.5,-0.5) rectangle (axis cs:14.5,1.5);
  \draw[kldim,dashed] (axis cs:0.5,0.5)--(axis cs:14.5,0.5) node[right,text=kldim]{candidate $a$};
  \node[text=kldim,anchor=west] at (axis cs:14.5,1.5) {$a+\varepsilon$};
  \node[text=kldim,anchor=west] at (axis cs:14.5,-0.5) {$a-\varepsilon$};
  \addplot[only marks,mark=*,mark size=1.4pt,klhi] coordinates {(2,1)(4,1)(6,1)(8,1)(10,1)(12,1)(14,1)};
  \addplot[only marks,mark=o,mark size=1.6pt,klink] coordinates {(1,-1)(3,-1)(5,-1)(7,-1)(9,-1)(11,-1)(13,-1)};
\end{axis}
\end{tikzpicture}
\end{center}
\plotcaption{Wherever you centre a band of half-width $\varepsilon=1$, it holds one of the two values,
never both. The other value (open dots here) escapes it infinitely often.}

\subsection{Generalize it}
\begin{itemize}
  \item \textbf{Two subsequences with different limits:} the sequence diverges, with
        $\varepsilon$ = half the distance between the limits.
  \item \textbf{No exact values needed:} if $a_n\ge c$ infinitely often and $a_n\le d$ infinitely
        often with $c>d$, the same proof works with $\varepsilon=(c-d)/2$.
  \item \textbf{Unbounded implies divergent:} the contrapositive of Problem 01, a second route to
        divergence.
\end{itemize}

\subsection{Same structure, elsewhere}
Problem 01 used convergence; this one proved its negation. In both, the move was choosing
$\varepsilon$ yourself. The uniqueness of limits runs on the same choice: two limits $a\neq b$, and
$\varepsilon=|a-b|/2$ makes the sequence sit in two disjoint bands at once.

\why{Why half the gap?}{A band of half-width $\varepsilon$ is $2\varepsilon$ wide. If $2\varepsilon$
is at most the distance between the two values, no band of that size holds both, wherever you
centre it. So whatever limit someone proposes, one value sits outside its band, and the sequence
returns to that value infinitely often.}

\why{Why does the order of the quantifiers matter?}{Swapped, $\varepsilon$ may depend on $N$ and
shrink as $N$ grows. Then $1/n$, which converges to $0$, satisfies the statement too: for each $N$
take $n=N$ and $\varepsilon=1/N$. A statement that convergent sequences also satisfy cannot prove
divergence.}

\subsection{The solution, idea first}
\emph{Idea:} every candidate is at least $1$ away from $1$ or from $-1$, and both are visited
infinitely often.
\begin{proof}
Let $a\in\mathbb{R}$. Since $2=|1-(-1)|\le|1-a|+|a-(-1)|$, one of the values $v\in\{1,-1\}$ has
$|v-a|\ge1$. Set $\varepsilon=1$. For every $N$ there is an $n\ge N$ with $a_n=v$ (an even or an odd
one), and for it $|a_n-a|\ge\varepsilon$. So $a$ is not a limit of $(a_n)$, and since $a$ was
arbitrary, $(a_n)$ diverges.
\end{proof}

\begin{triggers}
\trigger{A convergence hypothesis, and you need a bound}{choose $\varepsilon$ yourself;
  $\varepsilon=1$ if any will do}
\trigger{A claim about all terms of a convergent sequence}{split: the tail by one fixed
  $\varepsilon$, the finite head by a maximum}
\trigger{Show that something does not converge}{negate mechanically; two values visited infinitely
  often; $\varepsilon$ = half the gap}
\end{triggers}

\begin{deck}
\card{Trigger}{A convergence hypothesis, and you need a bound.}{Choose $\varepsilon$ yourself
  ($\varepsilon=1$ will do), then a maximum over the finite head.}
\card{Why}{Why may you choose $\varepsilon$ when convergence is given?}{``For every
  $\varepsilon$'' is a promise; you cash it in at any value you like.}
\card{Def}{$(a_n)$ does \emph{not} converge to $a$: $\exists\,\varepsilon>0\ \forall N\
  \exists\,n\ge N:$ \dots}{$|a_n-a|\ge\varepsilon$. Order kept, quantifiers and inequality flipped.}
\end{deck}

\clearpage
\section*{Again, in a week}
Due 23 October 2026. No notes. Write on the grid, then bring it to your weekly review. The solutions
are not in this file.

\problem[Variant of Specimen 1] Let $a_n\to a$ with $a\neq0$ and $a_n\neq0$ for all $n$. Show that
there is a $c>0$ with $|a_n|\ge c$ for all $n$.
\answerspace{85mm}

\problem[Variant of Specimen 2] Show that $b_n=\sin(n\pi/2)$ does not converge.
\answerspace{80mm}

\end{document}
